\subsection{The actual nonaffine integral normalization and its native points}

The actual relative-normalization construction is Mathlib's
\texttt{AlgebraicGeometry/Normalization.lean}, tagged Stacks 035H.
We apply it to the ORIGINAL generic-field morphism, rather than
requiring an affine-normalization or chart comparison as an input.
Geometric unibranchness means the STANDARD structural local-ring
property at every actual stalk, as above. Strict-henselization
characterizations and nonintegral normalization remain separate.

\begin{theorem}[The genuine generic morphism meets the exact normalization prerequisites]
\label{scheme:actual-integral-generic-normalization-input}
\lean{EnvelopeProperty.IntegralSchemeNormalization.genericMap,
EnvelopeProperty.IntegralSchemeNormalization.genericMap_apply,
EnvelopeProperty.IntegralSchemeNormalization.genericMap_quasiCompact,
EnvelopeProperty.IntegralSchemeNormalization.genericMap_quasiSeparated,
EnvelopeProperty.IntegralSchemeNormalization.genericMap_preimage,
EnvelopeProperty.IntegralSchemeNormalization.genericSectionIso,
EnvelopeProperty.IntegralSchemeNormalization.genericSectionIso_coefficients}
For ANY integral scheme $X$, put $K=X.functionField$, its ACTUAL
generic stalk. The real morphism
$j:\operatorname{Spec}K\to X$ is quasi-compact and quasi-separated.
For EVERY nonempty affine open $U$, its genuine inverse image is
all of $\operatorname{Spec}K$, and the entire actual section ring
$\Gamma(\operatorname{Spec}K,j^{-1}U)$ is genuinely isomorphic to
$K$ preserving the ORIGINAL coefficient map $\Gamma(X,U)\to K$.
That map makes $K$ an actual fraction field of the SAME original
$\Gamma(X,U)$. No quasi-compactness or separatedness of $X$ is needed.
\leanok
\end{theorem}
\begin{proof}
Use the genuine scheme morphism from the stalk at the original
generic point. Its source is the actual Spec of the actual field
$K$, so it has ONE point. Every inverse image is finite, hence
compact; this DERIVES quasi-compactness of the morphism. The actual
stalk morphism is a preimmersion, hence a monomorphism, hence
separated and therefore quasi-separated. These are derived from
the genuine map, not supplied as additional geometric hypotheses.

The source's unique point maps to the original generic point.
Every nonempty open of an integral scheme contains that point,
so its genuine inverse image is the whole source. Restrict sections
through this ACTUAL open equality and use the true global-sections
isomorphism of Spec $K$. The existing entire stalk-morphism section
formula identifies the WHOLE composite with the ORIGINAL germ map
to the generic stalk. Mathlib's actual function-field theorem makes
this same $K$ a fraction field of $\Gamma(X,U)$, using that genuine
coefficient action. No abstract field identification is assumed.
\end{proof}

\begin{theorem}[Actual relative normalization gives whole original affine closure charts]
\label{scheme:actual-integral-normalization-affine-charts}
\uses{scheme:actual-integral-generic-normalization-input,algebra:actual-normalization-fraction-field-choice}
\lean{EnvelopeProperty.IntegralSchemeNormalization.normalization,
EnvelopeProperty.IntegralSchemeNormalization.projection,
EnvelopeProperty.IntegralSchemeNormalization.genericLift,
EnvelopeProperty.IntegralSchemeNormalization.genericLift_comp_projection,
EnvelopeProperty.IntegralSchemeNormalization.projection_isIntegralHom,
EnvelopeProperty.IntegralSchemeNormalization.affineClosureEquiv,
EnvelopeProperty.IntegralSchemeNormalization.affineClosureEquiv_algebraMap,
EnvelopeProperty.IntegralSchemeNormalization.chartMap,
EnvelopeProperty.IntegralSchemeNormalization.chartMap_isOpenImmersion,
EnvelopeProperty.IntegralSchemeNormalization.chartMap_comp_projection,
EnvelopeProperty.IntegralSchemeNormalization.chartMap_opensRange}
Define $N(X)$ to be the ACTUAL relative normalization of $X$ in
$j:\operatorname{Spec}K\to X$, and let $\pi:N(X)\to X$ be its
true projection. It is integral. For each nonempty affine $U$,
$A=\Gamma(X,U)$, the LITERAL $C_U=\overline A^{\operatorname{Frac}A}$
has a genuinely constructed open immersion
$\operatorname{Spec}C_U\to N(X)$ whose ENTIRE range is
$\pi^{-1}U$. Its whole projection triangle is the ORIGINAL
$\operatorname{Spec}C_U\to\operatorname{Spec}A\to X$.
No normalization, affine cover, ring or scheme comparison is supplied.
The genuine generic-field morphism factors through $N(X)$ by the actual
normalization construction, and the WHOLE projection composite is the
original generic-field morphism.
\leanok
\end{theorem}
\begin{proof}
The exact qcqs prerequisites in the first node permit Mathlib's
true relative normalization. Its construction already proves
integrality of the projection and gives an ACTUAL open cover by
Spec of the integral closures of $\Gamma(X,U)$ in
$\Gamma(\operatorname{Spec}K,j^{-1}U)$; the range is precisely
$\pi^{-1}U$ and the whole projection map is the true original
coefficient map followed by the original affine chart.

The preceding actual ambient section isomorphism is an algebra
isomorphism over the ORIGINAL $A$. It maps true integral elements
to true integral elements. Actual fraction-field uniqueness then
compares this entire closure with the LITERAL $C_U$ through the
actual $K$ fraction-field action. These algebra isomorphisms
preserve EVERY original $A$ element. Apply Spec to the genuine
whole comparison and compose with the original normalization
cover injection. Its actual isomorphism and open immersion give
an open immersion, with the SAME actual inverse-image range.
The entire coefficient law proves the full original projection
triangle; selected coordinates or prime supports do not suffice.
\end{proof}

\begin{theorem}[True original normalization projection is scheme-surjective]
\label{scheme:actual-integral-normalization-surjective}
\uses{scheme:actual-integral-normalization-affine-charts}
\lean{EnvelopeProperty.IntegralSchemeNormalization.projection_surjective}
The genuine $\pi:N(X)\to X$ is scheme-surjective for ANY integral
$X$, without finite normalization, Noetherianity, quasi-compactness
or a chosen normalization witness.
\leanok
\end{theorem}
\begin{proof}
Given an ORIGINAL point of $X$, choose a genuine affine neighborhood
containing it; this neighborhood is nonempty because it contains
that actual point. The true affine chart identifies that point
with a prime of the SAME $A=\Gamma(X,U)$. Its LITERAL normalization
inclusion into $C_U$ is injective and integral, so actual lying-over
supplies a prime above it. The whole chart projection triangle in
the preceding node maps the resulting actual $N(X)$ point back
to the ORIGINAL point. Thus every whole scheme point is attained,
without choosing a support-lift certificate.
\end{proof}

\begin{theorem}[Actual open immersions preserve standard structural stalk geometry]
\label{scheme:actual-geometric-unibranch-open-inheritance}
\uses{algebra:actual-structural-unibranch-ring-equivalence}
\lean{EnvelopeProperty.IsGeometricallyUnibranchScheme.of_isOpenImmersion}
For an ACTUAL open immersion $Y\to X$, standard geometric
unibranchness of every genuine $X$ stalk implies the same standard
property for every genuine $Y$ stalk. No integrality or separation
hypothesis is needed for this local structural inheritance.
\leanok
\end{theorem}
\begin{proof}
For every ORIGINAL point $y$, the actual stalk map
$\mathcal O_{X,f(y)}\to\mathcal O_{Y,y}$ is a genuine ring
isomorphism, proved from the open immersion. Apply the already
proved true RingEquiv invariance, retaining its genuine reduction,
normalization and complete original residue square. This constructs
the target standard class; no normalization or native point
property is included in the hypothesis or conclusion.
\end{proof}

\begin{lemma}[The genuine nonempty affine opens cover every scheme]
\label{scheme:actual-nonempty-affine-open-cover}
\lean{EnvelopeProperty.IntegralSchemeNormalization.NonemptyAffineOpen,
EnvelopeProperty.IntegralSchemeNormalization.iSup_nonemptyAffineOpen_eq_top}
For ANY original scheme $X$, let the index family consist of its actual
nonempty affine opens. Their supremum in the original lattice of opens
is the whole scheme. No integrality, quasi-compactness, separatedness,
geometric unibranchness or supplied cover is assumed.
\leanok
\end{lemma}
\begin{proof}
For every original point $x\in X$, the existing affine-neighborhood
existence theorem supplies an actual affine open $U$ containing $x$.
That SAME point witnesses that $U$ is nonempty, so $U$ belongs to the
stated index family, and $x$ belongs to its supremum. Thus the supremum
is the top open. If $X$ is empty, the pointwise assertion is vacuous
and gives the same equality. No finite subcover is claimed.
\end{proof}

\begin{theorem}[Structural stalk geometry DERIVES global normalization radiciality]
\label{scheme:actual-integral-geometric-unibranch-normalization-radicial}
\uses{scheme:actual-integral-normalization-affine-charts,scheme:actual-geometric-unibranch-open-inheritance,scheme:actual-structural-unibranch-stalk-normalization,scheme:actual-nonempty-affine-open-cover}
\lean{EnvelopeProperty.IntegralSchemeNormalization.projection_universallyInjective}
If the actual integral $X$ is standard geometrically unibranch
at every ACTUAL local stalk, its genuine normalization projection
$\pi$ is universally injective. No prime-fibre, residue comparison,
MaxSpec injection, affine chart or native bijection is supplied.
\leanok
\end{theorem}
\begin{proof}
The ORIGINAL nonempty affine opens form a true cover, since every
scheme point has an actual affine neighborhood containing it.
For each such $U$, its actual chart $\operatorname{Spec}A\to X$
is an open immersion. The preceding structural inheritance theorem
derives the standard property of its actual scheme stalks. The
already proved integral-affine theorem therefore gives genuine
universal injectivity of the LITERAL $C_U$ projection.

The constructed open immersion $\operatorname{Spec}C_U\to N(X)$
and the true inclusion of $\pi^{-1}U$ have the SAME whole range.
Actual open-immersion uniqueness supplies a true isomorphism
between their source schemes. Its full projection identity, and
the original $\operatorname{Spec}A\simeq U$, identify the actual
restriction of $\pi$ with the SAME literal normalization map.
Transport its genuine algebraic universal injectivity through
those ACTUAL source and target isomorphisms. Universal injectivity
is genuinely Zariski-local on the target, so the DERIVED actual
nonempty affine cover proves it globally. No global or local
normalization comparison appears as an input.
\end{proof}

\begin{theorem}[The genuine nonaffine normalization preserves all native points]
\label{scheme:actual-integral-geometric-unibranch-normalization-native}
\uses{scheme:actual-integral-normalization-surjective,scheme:actual-integral-geometric-unibranch-normalization-radicial,scheme:integral-surjective-native-homeomorphism,scheme:relative-native-universal-homeomorphism}
\lean{EnvelopeProperty.IntegralSchemeNormalization.nativeHomeomorph,
EnvelopeProperty.IntegralSchemeNormalization.nativeHomeomorph_apply,
EnvelopeProperty.IntegralSchemeNormalization.relativeNativeHomeomorph,
EnvelopeProperty.IntegralSchemeNormalization.relativeNativeHomeomorph_apply}
For every integral scheme $X$ satisfying the STANDARD actual-stalk
geometric property, the independently defined native map
$N(X)^{\mathrm{val}}\to X^{\mathrm{val}}$ is a genuine homeomorphism.
It is EXACTLY the whole original projection on native points and
residue valuations. It likewise gives the exact whole homeomorphism
above ANY prescribed coefficient seminorm for ANY actual common
structure map $X\to\operatorname{Spec}R$.
\leanok
\end{theorem}
\begin{proof}
The actual projection is integral by construction, scheme-surjective
by true affine lying-over, and universally injective by the derived
actual stalk argument. The established genuine native invariance
producer therefore applies: true integral closedness and lifting,
plus actual purely inseparable residue valuation uniqueness, give
the native homeomorphism retaining the ORIGINAL whole map. Exact
basepoint composition restricts this SAME map above any common
prescribed coefficient point. No bounded base, field, completeness,
finite normalization, Noetherianity, quasi-compactness or separatedness
of $X$ is assumed. This is integral normalization, not a construction
for arbitrary nonintegral schemes, an equivalence with strict
henselization definitions, or an analytic structure-sheaf theorem.
\end{proof}
